\chapter{Planteamiento del problema}

\section{Planteamiento}

El avance de la robótica móvil ha impulsado el desarrollo de sistemas capaces de desplazarse de manera eficiente en entornos complejos, irregulares o no estructurados. Dentro de este campo, los robots cuadrúpedos se han consolidado como una alternativa relevante gracias a su capacidad para mantener estabilidad dinámica, adaptarse a variaciones del terreno y ejecutar patrones de marcha versátiles, lo que los convierte en plataformas idóneas para aplicaciones de exploración, inspección, rescate y automatización en ambientes donde otros robots presentan limitaciones.

En el ámbito académico, estos sistemas representan una herramienta fundamental para el estudio de la locomoción robótica, el modelado dinámico y el diseño de estrategias avanzadas de control. Su incorporación en espacios universitarios favorece el desarrollo de competencias investigativas y permite el análisis experimental de fenómenos asociados al movimiento, la estabilidad y la interacción mecánica con el entorno.

En este contexto, el laboratorio de robótica de la Universidad Santo Tomás incorporó en 2021 una plataforma cuadrúpeda adquirida a la empresa Nova Spot Micro, con el objetivo de fortalecer los procesos de formación e investigación en locomoción robótica. La plataforma fue concebida como un recurso didáctico para el desarrollo de proyectos orientados al análisis, modelado y control de sistemas dinámicos, y su adquisición se integró a los procesos institucionales de gestión de activos, permitiendo su registro y depreciación dentro del inventario académico.

Sin embargo, a pesar de su potencial, la plataforma no ha podido ser utilizada operativamente debido a múltiples limitaciones técnicas identificadas. En primer lugar, carece de un sistema electrónico funcional para el accionamiento de su estructura mecánica, lo que impide cualquier prueba de locomoción. A esto se suma la ausencia de un sistema adecuado de sensorización y potencia, restricciones que imposibilitan la validación de estrategias de control o la ejecución de experimentos de estabilidad y marcha.

Asimismo, no se dispone de un modelo matemático que describa su cinemática y dinámica ni de una estrategia de control que permita generar patrones de marcha estructurados, lo que limita la formulación de trayectorias articulares coherentes. Adicionalmente, la falta de documentación técnica detallada sobre sus características mecánicas, estructurales y geométricas dificulta su caracterización y correcta implementación.

Como consecuencia, la plataforma permanece inactiva desde su adquisición, reduciendo su impacto académico y restringiendo las posibilidades de investigación en locomoción cuadrúpeda dentro del laboratorio. Por ello, se hace necesario desarrollar un proceso integral que incluya la caracterización mecánica, el modelado cinemático así como un modelo dinámico, la implementación de un sistema electrónico funcional y el diseño de estrategias de control que permitan poner en operación la plataforma y cumplir con su propósito inicial.

\section{Pregunta problema}

¿En qué medida una estrategia de aprendizaje por refuerzo, utilizada como capa
correctiva de una marcha nominal, mejora la estabilidad y la repetibilidad de
la locomoción tipo paso y gateo del robot cuadrúpedo Nova Spot Micro en una
superficie plana y bajo condiciones controladas de laboratorio?
